% !TEX encoding = UTF-8 Unicode
% !TEX root =  ../Main.tex

% Rechtsquellen müssen in laut Richtlinien im Glossar eingetragen werden, nicht im Literaturverzeichnis. Damit diese nicht ausgeschrieben werden, wird der Befehl glsunset genutzt
\glsunset{dsgvozitat}
\glsunset{aeuv}
\glsunset{bdsgzitat}


\chapter{Theoretischer Rahmen}\label{cha:Theoretischer_Rahmen}

\section{Anforderungen an Datenschutz und Datensicherheit in Deutschland}\label{sec:Anforderungen}

Um eine Anwendung zu entwickeln, die personenbezogene Daten verarbeitet ist es notwendig, sich vorab mit den gesetzlichen Anforderungen zum Datenschutz auseinanderzusetzen. Die wichtigste Regularie dazu ist die \gls{dsgvo}. Es handelt sich hierbei um eine Verordnung der Europäischen Union und ist deshalb allgemein gültig, in all ihren Teilen verbindlich und gilt unmittelbar in jedem Mitgleidsstaat der Europäischen Union. [Vgl. §288 \gls{aeuv}]


Die DSGVO bestimmt Grundsätze und Vorschriften, um natürliche Personen bei der Verarbeitung ihrer personenbezogenen Daten zu schützen. Dadurch soll die Einhaltung ihrer Grundrechte sowie Grundfreiheiten gewährleistet werden. [Vgl. §1 \gls{dsgvozitat}] Die Verordnung gilt, mit wenigen Ausnahmen, für alle personenbezogenen Daten, die verarbeitet und in einem Dateisystem gespeichert werden. Die Verordnung gilt unabhängig davon, ob die Verarbeitung automatisiert oder nichtautomatisiert stattfindet. [Vgl. §2 \gls{dsgvozitat}] An dieser Stelle ist es notwendig, die Definitionen für "personenbezogene Daten" sowie "Dateisystem" genauer anzuschauen. Nach §4 der \gls{dsgvozitat} sind \glqq{}„personenbezogene Daten“ alle Informationen, die sich auf eine identifizierte oder identifizierbare natürliche Person (im Folgenden „betroffene Person“) beziehen; als identifizierbar wird eine natürliche Person angesehen, die direkt oder indirekt, insbesondere mittels Zuordnung zu einer Kennung wie einem Namen, zu einer Kennnummer, zu Standortdaten, zu einer Online-Kennung oder zu einem oder mehreren besonderen Merkmalen, die Ausdruck der physischen, physiologischen, genetischen, psychischen, wirtschaftlichen, kulturellen oder sozialen Identität dieser natürlichen Person sind, identifiziert werden kann;\grqq{}. Werden also Daten von Studierenden wie der volle Name und die Adresse durch den \gls{asta} verarbeitet, handelt es sich um personenbezogene Daten und die \gls{dsgvo} gilt. Ein Dateisystem ist nach der Definition der \gls{dsgvo} \glqq{}jede strukturierte Sammlung personenbezogener Daten, die nach bestimmten Kriterien zugänglich sind, unabhängig davon, ob diese Sammlung zentral, dezentral oder nach funktionalen oder geografischen Gesichtspunkten geordnet geführt wird;\grqq{}. Das bedeutet, dass jegliche Ablage der personenbezogenen Daten, sei es in einer Datenbank auf einem Server oder auf Papier in verschiedenen Ordnern, als Dateisystem zählt und auch hier die \gls{dsgvo} gilt.


Die Grundsätze der \gls{dsgvo} besagen, dass Daten nur auf eine rechtmäßige und für die betroffene Person nachvollziehbare Weise verarbeitet werden dürfen. Dabei muss darauf geachtet werden, dass nur die maximal notwendige Anzahl an Daten ausschließlich für den festgelegten Zweck verwendet wird. Die Daten müssen sachlich korrekt sein und dürfen nur für einen bestimmten Zeitraum vom Verwendenden gespeichert werden. Die Verarbeitung muss durch technische und organisatorische Maßnahmen geschützt sein. Der/die Verantwortliche muss dafür sorgen, dass diese Grundsätze eingehalten werden und diese Einhaltung auch nachweisen kann. [Vgl. §5 \gls{dsgvozitat}] Eine rechtmäßige Verarbeitung besteht beispielsweise dann, wenn die Person, deren personenbezogene Daten verarbeitet werden, dieser Verarbeitung zugestimmt hat. [Vgl. §6 Abs. 1 lit.~a \gls{dsgvozitat}] Die/der Verantwortliche muss eine solche Einwilligung auch nachweisen können. Die Einwilligung kann zudem jederzeit widerrufen werden.~[Vgl. §7 \gls{dsgvozitat}] Ein weiteres Szenario ist, wenn die Verarbeitung der Daten zur Erfüllung eines Vertrags notwendig ist, dessen Vertragspartei die betroffene Person ist. [Vgl. §6 Abs. 1 lit.~b \gls{dsgvozitat}] Weitere rechtsmäßige Szenarien sind in §6 \gls{dsgvozitat} beschrieben.  

Sobald die personenbezogenen Daten erhoben werden, ist die/der Verantwortliche dazu verpflichtet, der betroffenen Person den Namen und die Kontaktdaten der/des Verantwortlichen, den Verwendungszweck sowie die Rechtsgrundlage nach Artikel 6 zu nennen. [Vgl. §13 Abs. 1 \gls{dsgvozitat}] Der/die Verantwortliche muss zudem weitere Informationen zur Verfügung stellen, die für eine faire und transparente Verarbeitung notwendig sind. [Vgl. §13 Abs. 2 \gls{dsgvozitat}]

Neben den rechtlichen Grundlagen und Informationspflichten gibt es in der \gls{dsgvo} Regelungen dazu, wie die datenschutzkonforme Verarbeitung aussehen muss. Dabei handelt es sich um technische als auch organisatorische Maßnahmen. Dazu gehören Maßnahmen wie die Pseudonymisierung und Verschlüsselung der personnenbezogenen Daten, die Sicherstellung von Vertraulichkeit, Integrität, Verfügbarkeit und Belastbarkeit von genutzten Systemen und Diensten bei der Verarbeitung, die Möglichkeit Daten bei einem physischen oder technischen Vorfall wiederherstellen zu können und eine regelmäßige Überprüfung der technsichen und organisatorischen Maßnahmen zur Gewährleistung der Sicherheit der Daten. [Vgl. Art. 32 Abs. 1 \gls{dsgvozitat}]

Bei der Verletzung des Schutzes der personenbezogenen Daten muss, sofern ein Risiko für die Rechte und Feiheiten der Person vorliegt, die zuständige Aufsichtsbehörde benachrichtigt werden. [Vgl. Art. 33 \gls{dsgvozitat}]

Neben der \gls{dsgvo} auf europäischer Ebene gibt es das Pendant auf nationaler Ebene: das \gls{bdsg}. Es ergänzt diese Verordnung für Deutschland dort, wo Gestaltungspielraum gelassen wird. Das ist in §22 \gls{bdsgzitat} der Fall. In diesem werden besondere Schutzmaßnahmen für die Verarbeitung besonderer Kategorien personenbezogener Daten beschrieben. Darunter fallen Informationen zur rassischen und ethnischen Herkunft, politischen Meinung, religiosen und weltanschaulichen Überzeugung oder Gewerkschaftszugehörigkeit. Des weiteren fallen darunter auch genetische und biometrische Daten zur eindeutigen Identifizierung, Gesundheitsdaten, Daten zum Sexualleben oder zur sexuellen Orientierung. [Vgl. Art. 9 Abs. 1 \gls{dsgvozitat}]

Die \gls{dsgvo} steht über dem nationalen Gesetz, das bedeutet, Vorschriften des \gls{bdsg} finden keine Anwendung wenn in der \gls{dsgvo} etwas gegentieliges steht. [Vgl. Art. 1 Abs. 55 \gls{bdsgzitat}] Dieser Artikel ist besonders interessant, da es bereits Klagen gibt, die eine Nichteinhaltung dieses Gesetzes durch das \gls{bdsg} bemängeln. \VglZitat{eughUrteilVom300320242023}{} Das aktuelle \gls{bdsg} ist seit dem 25. Mai 2018 in Kraft, aktuell gibt es jedoch einen am 07.02.2024 verabschiedeten Regierungsentwurf für ein Gesetz zur Änderung des \gls{bdsg}, welches zum aktuellen Zeitpunkt noch nicht verkündet wurde. \VglZitat{WebsiteBMIDurchsetzungDatenschutzrecht}{2} Aus diesem Grund wird das \gls{bdsg} in dieser Arbeit nicht vollumfänglich analysiert. 



\section{Prozessbegriff}\label{sec:Prozess}